\section{Математическое описание}

\subsection{Граф и его представление}

Граф~--- пара $G = (V, E)$, где $V$~--- конечное непустое множество вершин,
$E \subseteq V \times V$~--- множество рёбер. Граф называется ориентированным,
если каждое ребро~--- упорядоченная пара, и неориентированным в противном
случае. Ациклический граф не содержит циклов; связный ациклический
неориентированный граф является деревом, ориентированный~--- DAG
(directed acyclic graph).

Взвешенный граф задаёт функцию весов $w : E \to \mathbb{R}$, ставящую
каждому ребру вещественное число. В данной работе веса генерируются
по распределению Рэлея--Райса с параметрами, задаваемыми пользователем.

\subsection{Распределение Рэлея--Райса}

Согласно справочнику Вадзинского (п.~3.16), обобщённое распределение Рэлея
(распределение Рэлея--Райса) задаётся плотностью вероятности:
\[
    f(z) = \frac{z}{a^2}\,
           e^{-(z^2 + h^2)/(2a^2)}\,
           I_0\!\left(\frac{zh}{a^2}\right),
    \quad z \geq 0,
\]
где $a > 0$~--- параметр масштаба, $h \geq 0$~--- параметр формы,
$I_0$~--- модифицированная функция Бесселя нулевого порядка.

Для генерации случайных чисел справочник приводит следующий алгоритм, если $u_1, u_2 \sim \mathcal{N}(0,1)$~--- независимые
стандартные нормальные значения, то
\[
    z = \sqrt{(h + a u_1)^2 + (a u_2)^2}.
\]


\subsection{Генерация связного ациклического графа по степеням вершин}

Пользователь задаёт число вершин~$n$ и параметры распределения $a$, $h$.
Число рёбер~$m$ не вводится напрямую.

\textbf{Шаг~1.} Для каждой вершины $i \in \{0, \ldots, n-1\}$ независимо
сэмплируется значение по формуле Вадзинского:
\[
    x_i = \sqrt{(h + a u_{i,1})^2 + (a u_{i,2})^2},
    \quad u_{i,1}, u_{i,2} \sim \mathcal{N}(0,1).
\]
Вектор $\mathbf{x} = (x_0, \ldots, x_{n-1})$ трактуется как ненормированные
веса вершин. Для перевода в вектор вероятностей применяется функция softmax
со стабилизацией :
\[
    \mathrm{softmax}(\mathbf{x})_i = p_i = \frac{e^{x_i - x_{\max}}}{\sum_{j=0}^{n-1} e^{x_j - x_{\max}}},
    \qquad \sum_{i=0}^{n-1} p_i = 1.
\]

\textbf{Шаг~2.} Итоговое число рёбер~$m$ сэмплируется из допустимого
диапазона $[n-1,\; n(n-1)/2]$, где $n-1$~--- минимум (остовное дерево),
$n(n-1)/2$~--- максимум (полный граф).

Для каждой вершины~$i$ определяется максимально допустимая степень
в топологическом порядке:
\[
    \mathit{capacity}_i = n - 1 - i.
\]
Взвешенная доля вершины~$i$: $q_i = p_i \cdot \mathit{capacity}_i$.
Вектор взвешенных долей масштабируется к сумме~$m$:
\[
    q_i \leftarrow q_i \cdot \frac{m}{\displaystyle\sum_j q_j}.
\]
Целевая степень вершины: $d_i = \lfloor q_i \rfloor$. Если после округления
$\sum_i d_i < m$, недостающие единицы последовательно прибавляются
к вершинам с наибольшей дробной частью $q_i - \lfloor q_i \rfloor$
при условии $d_i < \mathit{capacity}_i$.

\textbf{Шаг~3.} Граф строится в два прохода. В первом формируется
остовное дерево: для каждой вершины $i \geq 1$ предок выбирается
равномерно случайно из множества вершин $j < i$, для которых
остаток целевой степени $d_j > 0$; если таких вершин нет~---
равномерно из $\{0, \ldots, i-1\}$. После добавления ребра $(j, i)$
остаток целевой степени предка уменьшается: $d_j \leftarrow d_j - 1$.
Построение по индексам $j < i$ гарантирует связность и ацикличность.

Во втором проходе для каждой вершины~$i$ с $d_i > 0$ формируется
множество допустимых соседей $\{j \mid j > i,\; (i,j) \notin E\}$,
из которого случайно без возвращения выбираются
$\min(d_i,\,|\text{допустимых}|)$ вершин и добавляются соответствующие рёбра.

Веса рёбер сэмплируются по той же формуле Вадзинского.

\subsection{Эксцентриситет, центр и диаметральные вершины}

Для связного невзвешенного графа $G = (V, E)$ определим кратчайшее расстояние
$d(u, v)$ как суммарное расстояние минимального пути между~$u$ и~$v$.
Эксцентриситет вершины~--- максимальное расстояние до остальных вершин:
\[
    \varepsilon(v) = \max_{u \in V} d(v, u).
\]
На основе эксцентриситетов вводятся глобальные характеристики графа:
\[
    \mathrm{rad}(G)  = \min_{v \in V} \varepsilon(v), \qquad
    \mathrm{diam}(G) = \max_{v \in V} \varepsilon(v).
\]
Радиус графа --- множество вершин с минимальным эксцентриситетом:
\[
C(G) = \{v \mid \varepsilon(v) = \mathrm{rad}(G)\}
\]
Диаметральные вершины --- с максимальным:
\[
D(G) = \{v \mid \varepsilon(v) = \mathrm{diam}(G)\}
\]
Расстояния вычисляются алгоритмом обхода в ширину (\textit{BFS}), запускаемым из каждой вершины, алгоритмическая
сложность линейная ($O(V + E))$ для одного прохода, а итоговая сложность квадратична $O(V(V+E))$.

\subsection{Метод Шимбелла}

Пусть $A$~--- матрица весов графа размера $n \times n$: элемент $a_{ij}$
равен весу ребра $(i, j)$, если оно существует, иначе считается
неопределённым (\texttt{nullopt} в реализации). Метод Шимбелла позволяет
найти длины кратчайших и наидлиннейших путей из ровно~$k$ рёбер
для всех пар вершин одновременно.

Метод основан на переопределении двух алгебраических операций.
Операция «умножения» двух величин $a$ и $b$ заменяется их алгебраической
суммой:
\[
    a \cdot b = b \cdot a \;\Rightarrow\; a + b = b + a, \qquad
    a \cdot 0 = 0 \cdot a = 0 \;\Rightarrow\; a + 0 = 0.
\]
Операция «сложения» заменяется выбором минимального или максимального
из двух элементов:
\[
    a + b = b + a = \min(\max)\{a, b\},
\]
при этом нулевые (неопределённые) элементы игнорируются~--- минимум
или максимум выбирается из ненулевых значений.

На основе этих операций определяются два тропических матричных умножения.
Для операции $\otimes_{\min}$:
\[
    (A \otimes_{\min} B)_{ij}
    = \min_{k}\,\bigl(a_{ik} + b_{kj}\bigr),
\]
и аналогично $\otimes_{\max}$ с заменой $\min$ на $\max$. В общем виде
элемент матрицы $A^{(2)}$, полученной возведением $A$ в степень~2,
вычисляется как:
\[
    \omega_{ij}^{(2)} = \min(\max)
    \Bigl\{
        \omega_{i1}^{(1)} + \omega_{1j}^{(1)},\;
        \omega_{i2}^{(1)} + \omega_{2j}^{(1)},\;
        \ldots,\;
        \omega_{in}^{(1)} + \omega_{nj}^{(1)}
    \Bigr\}.
\]
Матрица $A^{(k)}$, вычисленная последовательным умножением начиная
с исходной матрицы $A$, содержит длины путей из ровно
$k$~рёбер для всех пар вершин. Временная сложность: $O(k \cdot n^3)$,
пространственная: $O(n^2)$.

\subsection{Подсчёт маршрутов}

Маршрут от вершины~$s$ до вершины~$t$~--- простой путь, не повторяющий
вершины. Все маршруты находятся методом обхода с возвратом (backtracking):
рекурсивный DFS с маской посещённых вершин. При достижении целевой вершины
текущий путь копируется в результат; при выходе из рекурсии вершина
снимается с отметки, что обеспечивает полный перебор всех вариантов.
Сложность в худшем случае~--- $O(n!)$, пространственная сложность $O(V)$.

\subsection{Обход рёбер в ширину}

Обход в ширину (\textit{BFS}) из вершины~$s$ посещает все вершины
графа. В основе алогитма лежит очередь~$Q$ и множество посещённых вершин~$\mathit{visited}$.
На каждом шаге вершина~$u$ извлекается из головы очереди; все
инцидентные рёбра $(u, v)$ обходятся, и каждая непосещённая вершина $v$
помечается и добавляется в хвост~$Q$.

Алгоритмическая сложность линейная $O(V + E)$ и пространственная сложность $O(V)$. Результатом обхода алгоритма является
последовательность рёбер в порядке их первого посещения и вектор
расстояний $d[v]$~--- минимальное число рёбер от~$s$ до~$v$.

\subsection{Алгоритм Дейкстры}

Алгоритм Дейкстры находит кратчайшие пути от источника~$s$ до всех
вершин во взвешенном графе с неотрицательными весами рёбер.
Поддерживается вектор расстояний $d[v]$, инициализированный~$+\infty$
для всех вершин кроме~$s$: $d[s] = 0$. На каждом шаге из приоритетной
очереди (min-heap) извлекается вершина~$u$ с минимальным~$d[u]$;
для каждого ребра $(u, v)$ с весом~$w(u,v)$ выполняется релаксация:
\[
    d[v] \leftarrow \min\!\bigl(d[v],\; d[u] + w(u,v)\bigr).
\]
Если значение обновилось~--- пара $(d[v], v)$ добавляется в очередь,
а предшественник $p[v] \leftarrow u$ записывается для восстановления пути.

Кратчайший путь до целевой вершины~$t$ восстанавливается обратным
проходом по массиву~$p$: начиная с~$t$, следуем $t \to p[t] \to
p[p[t]] \to \cdots \to s$ и разворачиваем последовательность.

Алгоритмическая сложность на бинарной куче:
$O((V + E)\log{V})$, пространственная сложность $O(V)$.

\subsection{Сеть потоков}

Сеть потоков~--- четвёрка $S = (G, c, s, t)$, где $G = (V, E)$~---
ориентированный граф, $c : E \to \mathbb{R}_+$~--- функция пропускных
способностей, $s \in V$~--- исток, $t \in V$~--- сток.
Поток в сети~--- функция $f : E \to \mathbb{R}_+$, удовлетворяющая
двум условиям: ограничению пропускной способности $f(u,v) \leq c(u,v)$
для всех $(u,v) \in E$, и закону сохранения потока~--- для каждой
вершины $v \neq s, t$:
\[
    \sum_{u : (u,v) \in E} f(u,v) = \sum_{w : (v,w) \in E} f(v,w).
\]
Величина потока: $|f| = \sum_{v : (s,v) \in E} f(s,v)$.
Остаточная пропускная способность ребра:
$r(u,v) = c(u,v) - f(u,v)$.

\subsection{Алгоритм Форда--Фалкерсона}

\textbf{Теорема} (Л.\,Р.\,Форд, Д.\,Р.\,Фалкерсон, 1962).
\textit{Величина максимального потока $p_{\max}$ в сети $S$ совпадает
с минимальной пропускной способностью $r_{\min}$ её разрезов.}

Алгоритм итеративно находит увеличивающий путь от истока к стоку
в остаточной сети~--- путь, по которому остаточная пропускная
способность каждого ребра строго положительна. Величина
дополнительного потока по найденному пути:
\[
    \Delta f = \min_{(u,v) \in P} r(u,v).
\]
Для каждого ребра $(u,v)$ пути $P$ поток обновляется:
\[
    f(u,v) \leftarrow f(u,v) + \Delta f, \qquad
    f(v,u) \leftarrow f(v,u) - \Delta f.
\]
Алгоритм завершается, когда увеличивающий путь не существует.
В реализации поиск пути выполняется через BFS (алгоритм
Эдмондса--Карпа), что гарантирует сложность $O(|V| \cdot |E|^2)$.

\subsection{Поток минимальной стоимости}

Каждому ребру $(u,v)$ дополнительно сопоставляется стоимость
$w(u,v) \in \mathbb{R}$. Стоимость потока:
\[
    \mathrm{cost}(f) = \sum_{(u,v) \in E} f(u,v) \cdot w(u,v).
\]
Задача~--- найти поток заданной величины $f^* = \lfloor\tfrac{2}{3}
\cdot |f_{\max}|\rfloor$ минимальной стоимости.

Алгоритм последовательно находит кратчайший (по стоимости) путь
от истока к стоку в остаточной сети алгоритмом Беллмана--Форда,
допускающим отрицательные веса рёбер. На каждой итерации по
найденному пути пропускается максимально возможный поток,
не превышающий $f^* - f_{\text{current}}$. Алгоритм Беллмана--Форда
выполняет $|V|-1$ релаксаций всех рёбер; итоговая сложность
одной итерации: $O(|V| \cdot |E|)$.

Остаточные рёбра обратного направления имеют стоимость $-w(u,v)$~---
это стандартный приём, позволяющий «отменить» ранее пущенный поток
при нахождении более дешёвого пути.
